\documentclass[10pt,a3paper]{article}
\usepackage[margin=10mm]{geometry}
\usepackage[T1]{fontenc}
\usepackage[utf8]{inputenc}
\usepackage[french]{babel}
\usepackage{lmodern}
\usepackage{microtype}
\makeatletter\def\input@path{{../../../Style/}}\makeatother
\NeedsTeXFormat{LaTeX2e}
\ProvidesPackage{TLPosterCarteMentale}[2026/08/03 Posters et cartes mentales SI]

\RequirePackage{tikz}
\RequirePackage[most]{tcolorbox}
\RequirePackage{xcolor}
\RequirePackage{amsmath,amssymb}
\RequirePackage{graphicx}
\usetikzlibrary{positioning,arrows.meta,calc,shapes.geometric,mindmap,backgrounds,fit}

\definecolor{TLPosterBlue}{RGB}{24,86,141}
\definecolor{TLPosterLight}{RGB}{234,243,250}
\definecolor{TLPosterAccent}{RGB}{232,126,34}
\definecolor{TLPosterGreen}{RGB}{52,152,102}
\definecolor{TLPosterRed}{RGB}{192,57,43}
\definecolor{TLPosterGray}{RGB}{80,90,100}

\newcommand{\TLPosterTitre}[2]{%
  \begin{tcolorbox}[enhanced,boxrule=0pt,arc=0pt,colback=TLPosterBlue,left=6mm,right=6mm,top=4mm,bottom=4mm]
    {\color{white}\bfseries\LARGE #1}\hfill{\color{white}\large #2}
  \end{tcolorbox}%
}

\newtcolorbox{TLPosterBloc}[2][]{
  enhanced,breakable,
  colback=TLPosterLight,
  colframe=TLPosterBlue,
  colbacktitle=TLPosterBlue,
  coltitle=white,
  fonttitle=\bfseries,
  title={#2},
  boxrule=0.8pt,
  arc=2mm,
  left=3mm,right=3mm,top=2mm,bottom=2mm,
  #1
}

\newtcolorbox{TLPosterFormule}[1][]{
  enhanced,
  colback=white,
  colframe=TLPosterAccent,
  boxrule=1pt,
  arc=2mm,
  fontupper=\bfseries,
  halign=center,
  #1
}

\newcommand{\TLPosterMethode}[1]{%
  \begin{tcolorbox}[enhanced,colback=TLPosterGreen!8,colframe=TLPosterGreen,title=Méthode,fonttitle=\bfseries]
  #1
  \end{tcolorbox}%
}

\newcommand{\TLPosterAttention}[1]{%
  \begin{tcolorbox}[enhanced,colback=TLPosterRed!6,colframe=TLPosterRed,title=Point de vigilance,fonttitle=\bfseries]
  #1
  \end{tcolorbox}%
}

\tikzset{
  TLmindroot/.style={concept,concept color=TLPosterBlue,text=white,font=\bfseries\large,minimum size=34mm},
  TLmindlevel1/.style={level 1 concept/.append style={font=\bfseries,minimum size=23mm,sibling angle=45}},
  TLmindlevel2/.style={level 2 concept/.append style={font=\small,minimum size=17mm}},
  TLarrow/.style={-{Stealth[length=3mm]},thick,draw=TLPosterBlue},
  TLbox/.style={rounded corners=2mm,draw=TLPosterBlue,fill=TLPosterLight,align=center,inner sep=3mm}
}

\newenvironment{TLCarteMentale}[1]{%
  \begin{tikzpicture}[mindmap,grow cyclic,concept color=TLPosterBlue,every node/.style={concept},TLmindlevel1,TLmindlevel2]
  \node[TLmindroot]{#1}
}{;
  \end{tikzpicture}
}

\newcommand{\TLBranche}[3][]{child[concept color=#2]{node[#1]{#3}}}

\endinput

\pagestyle{empty}
\renewcommand{\familydefault}{\sfdefault}
\begin{document}
\TLPosterCourseHeader{4}{Cinématique}{Position, vitesse, accélération et composition des mouvements}
\vspace{2mm}
\begin{minipage}[t]{0.61\linewidth}
\begin{TLPosterSavoirs}
\begin{itemize}
  \item définir trajectoire, vecteur position, vitesse et accélération dans un référentiel donné
  \item maîtriser produit vectoriel, vecteur rotation et dérivation d'un vecteur dans une base mobile
  \item utiliser le torseur cinématique et le champ des vitesses d'un solide indéformable
  \item composer vecteurs rotation, vecteurs vitesse et torseurs cinématiques entre repères
  \item connaître translation, rotation autour d'un axe, roulement sans glissement et loi de vitesse trapézoïdale
\end{itemize}
\end{TLPosterSavoirs}
\end{minipage}\hfill
\begin{minipage}[t]{0.36\linewidth}
\TLPosterKey{Deux méthodes pour $\vec V$}{1. Dériver le vecteur position dans le bon repère\\[1mm]2. Composer les mouvements puis utiliser\\$\vec V_Q=\vec V_P+\vec\Omega\wedge\overrightarrow{PQ}$}
\end{minipage}
\vfill
\begin{minipage}[t]{0.49\linewidth}
\begin{TLPosterBloc}[Méthode 1 : dériver une position]{TLPosterBlue}
Écrire $\overrightarrow{OP}$ dans une base adaptée puis calculer
\[
\vec V_{P/0}=\left.\frac{d\overrightarrow{OP}}{dt}\right|_0.
\]
Si la base tourne, dériver aussi ses vecteurs : la dérivation dans un repère mobile ajoute le terme de transport lié à $\vec\Omega$.
\end{TLPosterBloc}
\begin{TLPosterBloc}[Accélération]{TLPosterPurple}
Le cours privilégie la méthode robuste : une fois la vitesse obtenue,
\[
\vec a_{P/0}=\left.\frac{d\vec V_{P/0}}{dt}\right|_0.
\]
Exprimer d'abord tous les termes dans une base cohérente, puis dériver les scalaires \textbf{et} les vecteurs de base mobiles.
\end{TLPosterBloc}
\end{minipage}\hfill
\begin{minipage}[t]{0.49\linewidth}
\begin{TLPosterBloc}[Méthode 2 : torseur et composition]{TLPosterGreen}
Identifier une chaîne de mouvements simple. Composer d'abord les rotations puis les vitesses au \textbf{même point}. Transporter ensuite la vitesse avec le champ des vitesses. Cette méthode est particulièrement efficace lorsqu'un point de liaison possède une vitesse connue ou nulle.
\end{TLPosterBloc}
\begin{TLPosterBloc}[Cas particuliers utiles]{TLPosterTeal}
Translation : $\vec\Omega=\vec0$, tous les points ont la même vitesse. Rotation autour d'un axe fixe : choisir un point de l'axe, de vitesse nulle. Roulement sans glissement : la vitesse relative au point de contact est nulle. Loi trapézoïdale : pente de $v(t)$ = accélération et aire sous $v(t)$ = déplacement.
\end{TLPosterBloc}
\end{minipage}
\vfill
\begin{TLPosterMethode}
\begin{enumerate}
  \item Nommer précisément le mouvement étudié : solide / référentiel et point concerné.
  \item Choisir le repère et la méthode la plus courte : dérivation directe ou composition + champ des vitesses.
  \item Si nécessaire, déterminer le vecteur rotation à partir du paramétrage angulaire.
  \item Calculer la vitesse en évitant de transporter plusieurs fois inutilement le torseur.
  \item Pour l'accélération, dériver la vitesse dans le référentiel demandé.
  \item Vérifier direction, unité, cas particuliers ($\dot q=0$, point sur axe, contact sans glissement) et cohérence du mouvement réel.
\end{enumerate}
\end{TLPosterMethode}
\vfill
\TLPosterRetenir{choisir d'abord le point et le référentiel ; pour la vitesse, dérivation ou champ des vitesses ; pour l'accélération, dériver la vitesse.}
\end{document}
